%% --------------------------------------------------------
\section{Гессіан та нормальні моди}
%% --------------------------------------------------------

Коливання молекули відбуваються поблизу рівноважної геометрії $\mathbf{R}_0$.
Для малих відхилень потенціальну енергію апроксимують квадратичною формою:
\[
    E(\mathbf{R}) \approx E(\mathbf{R}_0) +
    \frac{1}{2} \Delta\mathbf{R}^T \, \mathbf{H} \, \Delta\mathbf{R},
\]
де $\mathbf{H}$ --- \textbf{гесіан} (матриця других похідних):
\[
    H_{ij} = \left. \frac{\partial^2 E}{\partial R_i \partial R_j} \right|_{\mathbf{R}_0}.
\]
Для $N$ атомів гесіан має розмір $3N \times 3N$ і повністю описує гармонічні коливання.

%% --------------------------------------------------------
\subsection{Матриця силових констант та діагоналізація}
%% --------------------------------------------------------

Елементи $H_{ij}$ --- це \textbf{силові константи}:
діагональні --- жорсткість окремого руху,
недіагональні --- зв’язок між різними рухами
(наприклад, розтяг \ce{C-H} та згин \ce{H-C-O}).

У декартових чи внутрішніх координатах коливання \emph{пов’язані}.
\textbf{Діагоналізація гесіану} дає:
\begin{itemize}
    \item \textit{власні значення} $\omega_k^2$ --- квадрати частот;
    \item \textit{власні вектори} $\mathbf{L}_k$ --- \emph{нормальні моди} (незалежні коливання).
\end{itemize}

Це роблять \textbf{після оптимізації геометрії}, щоб:
\begin{itemize}
    \item підтвердити мінімум ($\omega_k^2 > 0$);
    \item обчислити \textbf{ZPE}: $E_{\text{ZPE}} = \frac{1}{2} h \sum \omega_k$;
    \item знайти внесок у термохімію ($H$, $S$, $G$);
    \item змоделювати ІЧ/Раман-спектри (піки при $\omega_k$).
\end{itemize}

Отже, з гесіану через нормальні моди отримують повну картину коливальних властивостей молекули.

%% --------------------------------------------------------
\subsection{Нормальні коливання}
%% --------------------------------------------------------

Для переходу від гесіану $\mathbf{H}$ у декартових координатах до фізично осмислених частот враховують маси атомів.
Вводять \textbf{мас-зважений гесіан}:
\[
    \mathbf{F} = \mathbf{M}^{-1/2} \mathbf{H} \mathbf{M}^{-1/2},
\]
де $\mathbf{M}$ --- діагональна матриця атомних мас.

Далі розв'язують задачу на власні значення:
\[
    \mathbf{F} \mathbf{Q} = \omega^2 \mathbf{Q},
\]
де $\omega$ --- коливальні частоти, а $\mathbf{Q}$ --- вектори нормальних координат.

Отримані власні значення $\omega_i^2$ відповідають:
\begin{itemize}
    \item $3N - 6$ коливальним модам для нелінійної молекули;
    \item $3N - 5$ --- для лінійної;
    \item 3 трансляційним і 3 (або 2) обертальним модам з $\omega = 0$.
\end{itemize}

\paragraph{Зв'язок координат $Q_k$ та $\mathbf{R}$}

Нормальні координати $Q_k$ є лінійними комбінаціями атомних зміщень $\Delta \mathbf{R}$:
\[
    Q_k = \sum_{i=1}^{3N} \sqrt{m_i}\, l_{ik}\, \Delta R_i
\]
або у матричній формі:
\[
    \mathbf{Q} = \mathbf{L}^T \mathbf{M}^{1/2} \Delta \mathbf{R},
\]
де:
\begin{itemize}
    \item $\mathbf{M}$ — діагональна матриця мас атомів,
    \item $\mathbf{L}$ — матриця власних векторів (нормальних мод),
    \item $\Delta \mathbf{R} = \mathbf{R} - \mathbf{R}_0$ — відхилення координат від рівноваги.
\end{itemize}

Звідси обернене перетворення:
\[
    \Delta \mathbf{R} = \mathbf{M}^{-1/2} \mathbf{L} \, Q_k.
\]

Саме цю формулу використовують у чисельному диференціюванні:
\[
    \mathbf{R}^{(\pm)} = \mathbf{R}_0 \pm h\, \mathbf{M}^{-1/2} \mathbf{L}_k,
\]
що зв’язує крок $h$ у просторі нормальних координат $Q_k$ із фізичними зміщеннями атомів $\mathbf{R}$.


%% --------------------------------------------------------
\subsubsection{Масштабування частот}
%% --------------------------------------------------------

Розраховані гармонічні частоти $\nu_{\text{гарм}}$ систематично завищені відносно експериментальних $\nu_{\text{експ}}$ через спрощення моделі потенціальної поверхні та обмеження методів:

\begin{itemize}
    \item \textbf{Гармонічне наближення:} реальний потенціал ангармонічний, тому енергія коливань і частоти менші.
    \item \textbf{Нестача електронної кореляції:} метод Хартрі–Фока завищує жорсткість зв’язків.
    \item \textbf{Обмежений базис:} малий базис перебільшує градієнти енергії.
\end{itemize}

{Масштабування $\lambda$ емпірично компенсує ангармонічність і методичні похибки. Емпіричне виправлення:
\[
    \nu_{\text{корект}} = \lambda \cdot \nu_{\text{гарм}}
\]

\textbf{Типові масштабувальні фактори:}

\begin{center}
    \begin{tblr}{
        colspec = {X[l,2] X[r] X[l,2]},
        row{1} = {font=\bfseries},
        hline{1,2,Z} = {solid},
            }
        Метод/базис         & $\lambda$ & Коментар               \\
        HF/6-31G*           & 0.8929    & Сильно завищує         \\
        B3LYP/6-31G*        & 0.9613    & Універсальний          \\
        B3LYP/6-311+G(2d,p) & 0.9679    & Кращий базис           \\
        MP2/6-31G*          & 0.9434    & Для точних розрахунків \\
    \end{tblr}
\end{center}

\textit{Джерело: NIST Computational Chemistry Comparison and Benchmark Database}



%% --------------------------------------------------------
\subsection{Обчислення гесіану}
%% --------------------------------------------------------

PySCF може обчислювати гесіан аналітично або чисельно. Програма в файлі \inputcode{h2o_hessian.py} обчислює матрицю Гесіану для молекули води (\ce{H2O}) двома методами: чисельним та аналітичним.

\begin{commentbox}
У програмному пакеті \texttt{PySCF} гесіан зберігається не як звичайна
двовимірна матриця розмірності $3N \times 3N$, а у вигляді чотиривимірного
тензора
\[
H_{AB,ij} = \frac{\partial^2 E}{\partial R_{A,i} \partial R_{B,j}},
\]
де $A,B$~— індекси атомів ($1,\dots,N$), а $i,j \in \{x,y,z\}$~—
напрямки декартових координат.
Відповідно, внутрішня структура масиву має вигляд
\[
H \in \mathbb{R}^{N \times N \times 3 \times 3}.
\]

Щоб отримати стандартну матричну форму $3N \times 3N$,
необхідно спершу переставити осі, розмістивши напрямки координат поруч із
відповідними атомами:
\[
H'_{Ai,Bj} = H_{AB,ij}.
\]
Це виконується операцією
\[
\texttt{H'} = \inlinecode{H.transpose(0,2,1,3).reshape(3N,3N)}.
\]
Така трансформація формує коректну блочну структуру з блоками $3\times3$,
яка відповідає фізичному гесіану молекули у декартових координатах.
\end{commentbox}

\subsubsection{Чисельний Гесіан}

Чисельний Гесіан обчислюється через скінченні різниці методом центральних різниць:

\begin{enumerate}
    \item \textbf{Зміщення координат:} Кожна координата $x_i$ кожного атома зміщується на величину $\pm h$ (крок диференціювання):
    \begin{equation}
        x_i^{(\pm)} = x_i \pm h
    \end{equation}
    де $h = 0.001$ Bohr (атомні одиниці).

    \item \textbf{Обчислення градієнтів:} Для кожного зміщення обчислюється градієнт енергії:
    \begin{equation}
        \nabla E(x_i^{(+)}) = \left( \frac{\partial E}{\partial x_1}, \frac{\partial E}{\partial x_2}, \ldots, \frac{\partial E}{\partial x_N} \right)_{x_i = x_i + h}
    \end{equation}
    \begin{equation}
        \nabla E(x_i^{(-)}) = \left( \frac{\partial E}{\partial x_1}, \frac{\partial E}{\partial x_2}, \ldots, \frac{\partial E}{\partial x_N} \right)_{x_i = x_i - h}
    \end{equation}

    \item \textbf{Обчислення Гесіану:} Елементи матриці Гесіану обчислюються як:
    \begin{equation}
        H_{ij} = \frac{\partial^2 E}{\partial x_i \partial x_j} = \frac{\nabla E_j(x_i^{(+)}) - \nabla E_j(x_i^{(-)})}{2h}
    \end{equation}
\end{enumerate}

\textbf{Формула центральних різниць:}
\begin{equation}
    \boxed{H_{ij} = \frac{g_j(x_i + h) - g_j(x_i - h)}{2h}}
\end{equation}
де $g_j = \frac{\partial E}{\partial x_j}$ --- компонента градієнта.

\subsubsection{Аналітичний Гесіан}

Аналітичний Гесіан обчислюється через точні другі похідні енергії, використовуючи модуль PySCF:

\begin{equation}
    H_{ij} = \frac{\partial^2 E}{\partial x_i \partial x_j}
\end{equation}

Метод використовує \inlinecode{pyscf.hessian.rhf.Hessian(mf).kernel()} для прямого обчислення всіх елементів матриці Гесіану без числового диференціювання.


\paragraph{Обчислювальна складність}

\begin{itemize}
    \item \textbf{Чисельний метод:} Потребує $2 \times 3N$ SCF розрахунків, де $N$ --- кількість атомів.
    \begin{itemize}
        \item Для H$_2$O ($N=3$): $2 \times 3 \times 3 = 18$ SCF розрахунків.
        \item Час виконання: $\sim 10 - 30$ секунд.
    \end{itemize}

    \item \textbf{Аналітичний метод:} Один розрахунок з аналітичними похідними.
    \begin{itemize}
        \item Час виконання: $\sim 1 - 5$ секунд.
        \item \textbf{Прискорення:} у $5 - 20$ разів швидше.
    \end{itemize}
\end{itemize}

\paragraph{Точність}

\begin{itemize}
    \item \textbf{Аналітичний метод:} Точніший, оскільки використовує точні похідні без похибок дискретизації

    \item \textbf{Чисельний метод:} Точність залежить від кроку $h$:
    \begin{equation}
        \text{Похибка} \sim O(h^2) + \frac{\epsilon}{h}
    \end{equation}
    де $\epsilon$ --- машинна точність

    \item \textbf{Типова різниця:} Максимальна абсолютна різниця $\sim 10^{-5}$--$10^{-7}$ Ha/Bohr$^2$.
\end{itemize}

\subsubsection{Коливальні частоти}

З матриці Гесіану можна отримати коливальні частоти молекули:

\begin{equation}
    \omega_i = \sqrt{\lambda_i} \times 5140.48,
\end{equation}
де $5140.48$~--- це числовий множник, який переводить частоти з атомних одиниць у~см$^{-1}$, а $\lambda_i$ --- власні значення масово-зваженого Гесіану:

\begin{equation}
    \tilde{H}_{ij} = \frac{H_{ij}}{\sqrt{m_i m_j}}
\end{equation}

\noindent\textbf{Очікувані частоти для H$_2$O:}
\begin{itemize}
    \item Перші 6 мод: $\approx 0$ см$^{-1}$ (трансляції та обертання)
    \item Симетричне валентне коливання: $\approx 3600$ см$^{-1}$
    \item Деформаційне коливання: $\approx 1600$ см$^{-1}$
    \item Асиметричне валентне коливання: $\approx 3800$ см$^{-1}$
\end{itemize}

Для практичних розрахунків \textbf{аналітичний метод є кращим вибором}, оскільки він значно швидший та точніший. Чисельний метод корисний для перевірки результатів або коли аналітичні похідні недоступні.
